\section{USB、HID 与音频接口的共同边界}

USB 是主机调度的分层总线。设备接入后先完成复位和枚举，主机读取描述符、分配地址、选择配置，再由驱动按 endpoint 类型提交传输。设备不能像 PCIe 主设备那样任意 DMA 系统内存，而是由主控制器按调度表发起总线事务，再由主控制器 DMA 主机缓冲区。

\begin{longtable}{L{0.18\linewidth}L{0.32\linewidth}L{0.4\linewidth}}
\toprule
传输类型 & 典型设备 & 优先保证 \\
\midrule
control & 枚举与配置 & 正确、可重试、双向短消息 \\
bulk & U 盘、打印机 & 正确性和剩余带宽利用率 \\
interrupt & 键盘、鼠标、触控 & 有界轮询间隔，不等同于物理中断线 \\
isochronous & 音频、摄像头 & 周期带宽和时延，允许按协议处理丢失 \\
\bottomrule
\end{longtable}

HID 报告把按键或指针状态编码成固定格式；音频接口还要处理生产者与消费者时钟略有偏差的问题，靠反馈端点、缓冲水位或采样率转换避免缓冲区逐渐溢出/耗尽。传感器接口也遵循同一思路：先确定采样时钟和数据格式，再定义缓冲、时间戳、丢样与错误报告。

\section{从物理事件到进程可见事件}

一次键盘按键会经历开关抖动、设备端扫描与消抖、USB 报告、主控制器 DMA、完成中断、内核输入子系统和用户态读取。每层都可能合并或延迟事件，因此“按键发生时刻”至少有物理接触、设备采样、主机收包和进程读取四个时间戳。

排错时按层观察：设备是否枚举，endpoint 是否配置，主控制器队列是否前进，DMA 缓冲是否更新，完成中断是否到达，驱动是否把报告交给上层。只在应用程序里反复重试，无法区分线缆、电源、协议、控制器和驱动故障。
